\chapter{项目概述与主张}

\section{项目定义}

Pebble 是一个面向复杂人际互动场景的 AI 情绪防御与沟通辅助系统。本文对该项目的核心定义是：系统以关系上下文为解释中心，把对话解码、关系陪练、稳态恢复与关系档案管理组织为同一条用户支持链路，使用户能够在高压沟通中先理解局势，再选择应对方式，并在后续互动中持续积累策略经验。

这一项目定义同时包含两个层次。

\begin{itemize}
  \item Pebble 是产品对象，关注用户在真实关系中的表达压力、误读风险与边界维护需求。
  \item Pebble 是系统对象，要求前端交互、分析服务、关系数据与反馈机制能够围绕同一组关系上下文协同工作。
\end{itemize}

后续章节会分别展开这两个层次，但本章先把项目主张收束为统一的问题陈述。

\section{目标与非目标}

Pebble 的直接目标有三项。

\begin{itemize}
  \item 帮助用户更准确地解释复杂对话中的潜台词、情绪信号与关系风险。
  \item 帮助用户把一次性的理解结果转化为可练习、可比较、可回顾的回应能力。
  \item 帮助用户在高压时刻获得最小可用的稳定支持，并把关系背景纳入后续分析，使系统逐步形成更细粒度的策略适配能力。
\end{itemize}

本项目的非目标同样需要明确。Pebble 不以心理诊断、医学治疗或人格定性为目标，也不尝试替代用户在真实关系中的判断责任。系统可以提供分析、建议与练习环境，但最终如何理解对方、是否采纳建议、是否调整关系边界，仍由用户自行决定。后续涉及画像、洞察或策略提示的章节，都应在这一边界内理解其作用。

\section{核心价值主张}

Pebble 的核心价值不在于把单句文本送入模型后返回一段解释，而在于把“关系背景”持续纳入分析过程。对于同一句话，来自母亲、伴侣、同事或上级的表达，往往意味着不同的互动历史、风险结构与回应策略。若系统忽略关系上下文，就只能输出抽象而平均化的建议；若系统能够维护关系档案与长期观察，它就有条件在解释与建议中引入更贴近用户现实处境的差异化判断。

这一价值主张进一步带来两个产品层面的意义。

\begin{itemize}
  \item Pebble 可以从一次性分析工具演进为持续性关系洞察系统。
  \item Pebble 可以把“理解”与“应对”连接起来，使用户不只得到一个分析结果，还能进入陪练、记录、复盘与下次适配的循环。
\end{itemize}

这一循环决定了项目后续架构、数据组织和页面设计的基本方向。

表~\ref{tab:pebble-project-claim} 对本章的项目主张做进一步压缩。它的作用是把 Pebble 的系统对象、直接目标与边界约束放在同一张表中，便于读者在进入背景章之前先形成稳定的比较框架。

\begin{table}[H]
  \centering
  \caption{Pebble 项目主张概览，突出系统对象、目标与边界的对应关系。}
  \label{tab:pebble-project-claim}
  \begin{tabularx}{\textwidth}{>{\raggedright\arraybackslash}p{2.6cm} >{\raggedright\arraybackslash}X >{\raggedright\arraybackslash}X}
    \toprule
    维度 & Pebble 的主张 & 本章给出的边界 \\
    \midrule
    系统对象 & 关系上下文驱动的 AI 情绪防御与沟通辅助系统 & 同时覆盖产品体验、分析服务与关系数据协同 \\
    直接目标 & 帮助用户理解高压对话、训练回应能力，并在需要时快速稳定状态 & 不把单次情绪识别视为最终目标，而强调连续支持链路 \\
    差异化来源 & 关系档案与长期观察可让同一句话获得更贴近情境的解释与建议 & 不承诺脱离关系背景即可给出高质量判断 \\
    非目标 & 系统提供分析、建议与练习环境 & 不承担心理诊断、医学治疗或替代用户判断的职责 \\
    报告定位 & 以后续章节展开产品结构、背景、系统设计与验证路径 & 当前章只立题，不展开接口、迁移历史和工程阻塞 \\
    \bottomrule
  \end{tabularx}
\end{table}

\section{报告主张与阅读位置}

本报告对 Pebble 的主张可以概括为三点。

\begin{itemize}
  \item Pebble 应被理解为关系上下文驱动的沟通支持系统，而非孤立的情绪识别工具。
  \item Pebble 的能力结构必须以解码、练习、稳态恢复与长期记忆的联动来描述，缺失其中任一环节都会削弱系统闭环。
  \item Pebble 的产品价值来自用户问题、关系数据与交互流程之间的协同设计，因此后续分析不能只停留在页面视觉或接口列表层面。
\end{itemize}

从全篇结构看，本章承担“立题”功能。第三章会进一步解释这一项目为何成立，并把背景、冲突与上下文展开为更完整的问题域分析；第四章至第八章会把术语、原则、流程、系统设计与验证路径依次展开；第九章与第十章则分别收束当前进展与风险边界。读者完成本章后，应已理解 Pebble 关注什么问题、其系统边界为何成立，以及后续各章将如何围绕这一主张逐步展开。

\section{本章小结}

本节读完后，读者应能说明 Pebble 的直接目标与非目标边界，并解释为什么系统价值不能只用单句情绪识别来理解。
